% Generated draft: requires mathematical and source review.
\documentclass[11pt]{article}
\usepackage[margin=25mm]{geometry}
\usepackage{fontspec}
\setmainfont{Latin Modern Roman}
\usepackage{amsmath,amssymb,mathtools,unicode-math}
\setmathfont{Latin Modern Math}
\usepackage{enumitem,needspace,xurl,hyperref}
\setlength{\parindent}{0pt}
\setlength{\parskip}{5pt}
\setlength{\emergencystretch}{3em}
\newcommand{\N}{\mathbb{N}}
\newcommand{\Z}{\mathbb{Z}}
\newcommand{\Q}{\mathbb{Q}}
\newcommand{\R}{\mathbb{R}}
\newcommand{\C}{\mathbb{C}}
\newcommand{\F}{\mathbb{F}}
\newcommand{\A}{\mathbb{A}}
\newcommand{\PP}{\mathbb{P}}
\newcommand{\E}{\mathbb{E}}
\newcommand{\eps}{\varepsilon}
\newcommand{\dd}{\,\mathrm{d}}
\begin{document}
\section*{UnsolvedMath ID: 30004659}
\textbf{Source problem code:} catalogue entry
\subsection{Definitions and mathematical statement}
Write \(\mathcal O[\Delta_Q]\) for the group algebra of the source's coefficient ring \(\mathcal O\) and group \(\Delta_Q\). Retain the source's notation \(R_{S\cup Q}\) for its deformation ring, \(\rho_{S\cup Q}\) for its associated Galois representation, and \(\bar\rho\) for the residual representation determining the maximal ideal \(\mathfrak m_{\bar\rho}\). Set
\[
A=\mathcal O[\Delta_Q],\qquad
R=R_{S\cup Q},\qquad
T=\left(\mathbb T_{S\cup Q,Q\text{-aug}}(K_0(Q)/K_1(Q))\right)_{\mathfrak m_{\bar\rho}}.
\]
Here the subscript denotes localization at the indicated maximal ideal. The notation \(Q\text{-aug}\), the level expression \(K_0(Q)/K_1(Q)\), and the intended \(A\)-algebra structures are retained from the source; their constructions are not supplied in the record.

For each source-intended index \(v\notin S\cup Q\), let
\[
\chi_v(X)=\operatorname{charpoly}\!\left(\rho_{S\cup Q}(\operatorname{Fr}_v);X\right)\in R[X],
\]
and let \(P_v(X)\in T[X]\) denote the source's Hecke polynomial, viewed in the localized algebra. These coefficient-ring assignments express the intended compatibility in the supplied statement; the record does not supply the hypotheses or constructions that establish them. For an \(A\)-algebra homomorphism \(\varphi:R\to T\), write \(\varphi_X:R[X]\to T[X]\) for its coefficientwise extension, defined by
\[
\varphi_X\!\left(\sum_i a_iX^i\right)=\sum_i\varphi(a_i)X^i.
\]

\textbf{Conjectural statement, relative to the intended source setup.} Does there exist an \(A\)-algebra homomorphism
\[
\varphi:R\longrightarrow T
\]
such that
\[
\varphi_X(\chi_v(X))=P_v(X)
\qquad\text{for every }v\notin S\cup Q?
\]

\textbf{Unresolved qualifications.} This is a source-relative formulation, not a fully specified standalone conjecture. The record does not identify the underlying global field or Galois group, the coefficient ring and residual representation hypotheses, the admissible choices of \(S\) and \(Q\), the deformation conditions defining \(R\), the construction of \(\Delta_Q\) and its actions, the meaning of the augmented Hecke algebra and its levels, or the definitions and normalizations of \(\operatorname{Fr}_v\) and \(P_v(X)\). It also does not specify the range of arithmetic setups over which the conjecture is intended to hold. No universal quantifier over unspecified setups is added here; these data must be recovered from the cited source before a standalone statement can be certified.
\subsection{Short English statement}
In the source's intended Taylor--Wiles setup, the conjecture asks for an \(\mathcal O[\Delta_Q]\)-algebra map from the deformation ring to the localized \(Q\)-augmented Hecke algebra that identifies every indicated Frobenius characteristic polynomial with its Hecke polynomial outside \(S\cup Q\). The supplied record leaves essential setup and normalization data unspecified, so this draft requires source consultation and human mathematical review.
\subsection{Sources}
\begin{itemize}
\item Local UnsolvedMath registry, supplied complete source record for permanent ID \(30004659\), exact source problem code \texttt{catalogue entry}, title: Taylor--Wiles Deformation Maps to Augmented Hecke Algebras. The record assigns proposed year \(2021\) and names no individual proposer.
\item Supplied source citation: \emph{Arbeitsgemeinschaft: Derived Galois Deformation Rings and Cohomology of Arithmetic Groups} (2022), Oberwolfach Reports, pp.\ 1001--1046, DOI \texttt{10.4171/owr/2021/18}. \url{https://doi.org/10.4171/owr/2021/18}. The background describes these reports as documenting mathematical workshops and their open-problem discussions. The literature-source label is ``Oberwolfach Report 18/2021.'' The supplied citation's year \(2022\) and the record's \(2021\) report/proposal labels are retained without reconciliation.
\item Supplied literature reference: Galatius--Venkatesh, ``derived Galois deformation rings.'' \url{https://arxiv.org/abs/1608.07236}.
\item Supplied literature assessment, checked exactly \texttt{2026-08-22}: the record labels the problem open and describes this Taylor--Wiles-level deformation-to-Hecke algebra map as a conjectural input. It states that derived patching develops consequences conditional on such inputs and that no general proof of this exact statement was found. This is supplied dated curation derived from the OWR material, not an independently verified literature assessment; time-sensitive status and cited claims require independent verification.
\end{itemize}
\end{document}
